% Chapter 11: Conjunctions: The Joining Words

\chapter{Conjunctions: The Joining Words}

\section{Pedagogical Purpose}
After learning about individual parts of speech, students are ready to understand how words and ideas are connected. Conjunctions are essential for building more complex sentences, improving sentence variety, and making writing and speaking more coherent and less choppy.

\begin{learningobjectives}
The learner can:
\begin{itemize}
    \item define what a conjunction is.
    \item identify conjunctions in sentences.
    \item use coordinating conjunctions (for, and, nor, but, or, yet, so) correctly.
    \item combine two simple sentences into one compound sentence using a coordinating conjunction.
    \item understand how conjunctions improve sentence flow and variety.
\end{itemize}
\end{learningobjectives}

\section{Prerequisite Check}
\begin{quickcheck}
Combine these pairs of sentences into one sentence, if possible.
\begin{enumerate}
    \item Rohan likes apples. Maya likes apples.
    \item The sun is shining. It is cold.
\end{enumerate}
\end{quickcheck}

\section{Language Experience (Let's Begin)}
Ms. Sharma is talking to Maya and Rohan about their plans for the weekend.

\textbf{Ms. Sharma:} "What are your plans for Saturday?"

\textbf{Rohan:} "I will play football. I will watch a movie."

\textbf{Maya:} "I will read a book. I will draw a picture."

\textbf{Ms. Sharma:} "Those are good plans! Rohan, can you tell me your plans in one sentence?"

\textbf{Rohan:} "I will play football \textbf{and} I will watch a movie."

\textbf{Ms. Sharma:} "Excellent! And Maya, what about you?"

\textbf{Maya:} "I will read a book \textbf{or} I will draw a picture."

\textbf{Ms. Sharma:} "Perfect! The words 'and' and 'or' are very useful. They help us join words or sentences together. These joining words are called \textbf{Conjunctions}."

\section{Concept Discovery (What is a Conjunction?)}
A \textbf{conjunction} is a word that joins words, phrases, or sentences together.

Conjunctions act like bridges, connecting different parts of a sentence.

\begin{examplebox}
    Rohan \textbf{and} Maya are friends.
    \begin{itemize}
        \item \texttt{and} joins the words \texttt{Rohan} and \texttt{Maya}.
    \end{itemize}
\end{examplebox}

\begin{examplebox}
    I like to read \textbf{but} I don't like to write.
    \begin{itemize}
        \item \texttt{but} joins two sentences.
    \end{itemize}
\end{examplebox}

\section{Concept Explanation (Coordinating Conjunctions)}
\textbf{Coordinating conjunctions} join words, phrases, or sentences that are of equal importance. They connect elements that are grammatically similar.

There are seven coordinating conjunctions, and you can remember them with the acronym \textbf{FANBOYS}:

\begin{itemize}
    \item \textbf{F}or (shows reason)
    \item \textbf{A}nd (adds information)
    \item \textbf{N}or (shows negative choice)
    \item \textbf{B}ut (shows contrast)
    \item \textbf{O}r (shows choice)
    \item \textbf{Y}et (shows contrast, similar to 'but')
    \item \textbf{S}o (shows result)
\end{itemize}

\section{Guided Practice (Let's Practise Together)}
\begin{activitybox}{Activity A: Choose the Right Conjunction}
\textbf{Ms. Sharma:} "Let's choose the best conjunction to connect these ideas. I'll do the first one."
\begin{enumerate}
    \item I wanted to buy the book, \_\_\_\_\_\_ I did not have enough money. (so / but)
    \item Maya is clever \_\_\_\_\_\_ she is also very kind. (and / or)
    \item Should I have the apple, \_\_\_\_\_\_ should I have the banana? (or / so)
    \item It was late, \_\_\_\_\_\_ we went home. (but / so)
\end{enumerate}
\end{activitybox}

\section{Summary}
\begin{summarybox}
In this chapter, we learned that \textbf{conjunctions} are joining words.
\begin{itemize}
    \item They connect words, phrases, and sentences.
    \item \textbf{Coordinating Conjunctions} join equal parts. You can remember them with \textbf{FANBOYS}.
    \item \textbf{F}or (shows a reason) - I was tired, for I had walked a long way.
    \item \textbf{A}nd (adds information) - I like apples and bananas.
    \item \textbf{N}or (shows a negative choice) - He does not play, nor does he study.
    \item \textbf{B}ut (shows contrast) - It is sunny, but it is cold.
    \item \textbf{O}r (shows a choice) - Do you want tea or coffee?
    \item \textbf{Y}et (shows contrast) - He is rich, yet he is unhappy.
    \item \textbf{S}o (shows a result) - It was raining, so we stayed inside.
\end{itemize}
Using conjunctions makes your writing flow better and helps you express more complex ideas in a single sentence.
\end{summarybox}